\BeleskeID{09_3-s010}
% element: 09_3-s010:e04
% element: 09_3-s010:e05
% slide-id: 09_3-s010
$C_L$ je ukupna kapacitivnost kojom je opterećen prvi invertor. Ona obuhvata ukupne kapacitivnosti oba drejna u prvom stepenu i oba gejta u drugom.

Kada ulaz poraste sa $0$ na $V_{DD}$, izlaz opada sa $V_{DD}$ na $0$. Promena napona između krajeva kapacitivnosti $C_{GDp1}$ zato je približno $2V_{DD}$. Količina naelektrisanja je $\Delta Q\approx C_{GDp1}\,2V_{DD}$. Da bi se isto naelektrisanje predstavilo kapacitivnošću prema nepokretnoj referenci na izlazu, gde promena iznosi samo $V_{DD}$, ekvivalentna kapacitivnost mora iznositi približno $2C_{GDp1}$. Drejnska kapacitivnost uključuje i $C_{DBp1}$. Isto razmišljanje važi za N tranzistor.

U oba invertora širine P kanala povećavamo, dok N širine ostaju iste. Faktor $\beta$ predstavlja odnos geometrijskih faktora. $C_{Dp1(\beta=1)}$ je kapacitivnost pre povećanja, pri $W_{n1}=W_{p1}$. U upotrebljenoj aproksimaciji ona se izjednačava sa $C_{Dn1}$, pa ukupan P doprinos raste približno linearno sa $\beta$.
